\documentclass{article}
\usepackage{enumitem}
\usepackage{neil}

%\setlength{\parindent}{0pt}

\title{MS Algebra Recitation 3}
\author{Neil Makur}
\date{}

\begin{document}

\maketitle

\section{Good Group Actions to Know 2: Electric Boogaloo}

\begin{defn*} Let $G\acts G$ by conjugation (i.e. $g\cdot h=ghg^{-1}$). The \textit{centralizer} of $h\in G$ is the stabilizer of $h$ under this action. That is,
$$Z_G(h)=G_h=\{g\in G:ghg^{-1}=h\}=\{g\in G:gh=hg\}.$$
\end{defn*}

\begin{defn*} Let $G\acts G$ by conjugation. The \textit{center} of $G$ is the kernel of the map $G\to S_G$. That is,
$$Z(G)=\{g\in G:ghg^{-1}=h~\forall h\in G\}=\{g\in G:gh=hg~\forall h\in G\}.$$
\end{defn*}

\begin{defn*} Let $G\acts\{H:H\leq G\}$ by conjugation (verify that this is a well-defined action!). The \textit{normalizer} of $H$ is the stabilizer of $H$ under this action. That is,
$$Z_G(H)=G_H=\{g\in G:gHg^{-1}=G\}.$$
The \textit{centralizer} of $H$ is the pointwise normalizer of $h$:
$$Z_G(H)=\{g\in G:ghg^{-1}=h~\forall h\in H\}=\bigcap_{h\in H}Z_G(h).$$
\end{defn*}

\noindent Because of the way we defined these, we immediately get the following:
\begin{itemize}
  \item Centralizers are subgroups
  \item Normalizers are subgroups
  \item The center is a \textit{normal} subgroup
\end{itemize}

Also straightforward from the definition is that $H\unlhd N_G(H)$ for any $H\leq G$. This ends up being pretty useful, since even if we don't have a normal subgroup, we can restrict and assume our subgroup is normal.

The following are exercises in symbol pushing:
\begin{itemize}
  \item $g\in Z_G(h)$ if and only if $h\in Z_G(g)$
  \item $Z(G)$ is the set of fixed points of the action $G\acts G$ by conjugation
  \item $Z_G(g)$ is the set of fixed points of $\varphi_g:h\mapsto ghg^{-1}$
  \item $Z(G)=\bigcap_{g\in G}Z_g(g)$
\end{itemize}

\section{Why Do We Care About Normal Subgroups?}

We first give an easy way to generate normal subgroups.

\begin{prop*} Let $G$ be a group and $\phi:G\to H$ be a homomorphism. Then $\ker\phi\unlhd G$.\end{prop*}
\begin{proof} Recall that $K\unlhd G$ if and only if $gKg^{-1}\subseteq K$ for all $g\in G$. Set $K=\ker\phi$ and take $k\in K$ and $g\in G$. Then,
$$\phi(gkg^{-1})=\phi(g)\phi(k)\phi(g)^{-1}=\phi(g)1_H\phi(g)^{-1}=\phi(g)\phi(g)^{-1}=1_H.$$
So $gKg^{-1}\subseteq K$ and $K\unlhd G$.
\end{proof}

To motivate the definition of a normal subgroup, we discuss how one would want to construct a quotient group. Suppose $G$ is a group and let $\sim$ be an equivalence class on $G$. Then we want to make $G/\sim$ into a group via $[g][h]=[gh]$. For this to be well-defined, it needs to be equivalent regardless of what representative we pick for each equivalence class. That is, if $g\sim g'$ and $h\sim h'$, then we need $gg'\sim hh'$. We'll consider a slight weakening of this property.

\begin{prop*} Let $G$ be a group and $\sim$ be an equivalence relation on $G$ with $g\sim g'\implies hg\sim hg'$ for all $g,g',h\in G$. Then $H:=[1]$ is a subgroup of $G$, and $[g]=gH$.\end{prop*}
\begin{proof} We first verify that $H\leq G$. Clearly $1\in H$ so $H$ is nonempty. If $g,g'\in H$ then we have $1\sim g$ and $1\sim g'$. Multiplying $1\sim g'$ on the left by $g$ gives $g\sim gg'$. Transitivity then gives $1\sim g\sim gg'$. Similarly, multiplying $1\sim g$ on the left by $g^{-1}$ gives $g^{-1}\sim 1$, and reflexivity gives $1\sim g^{-1}$. Thus, $H\leq G$.

Now, if $g\sim g'$ then $1\sim g^{-1}g'$ and so $g^{-1}g'\in H$. Thus, $gH=g'H$ and so the equivalence class of $g$ is contained in $gH$. Conversely, if $gH=g'H$ then $g^{-1}g'\in H$, so $g^{-1}g'\sim 1$ and $g'\sim g$. Thus, $[g]=gH$.
\end{proof}

Symmetrically, if $\sim$ satisfies $g\sim g'\implies gh\sim g'h$, then the equivalence classes are the right cosets of $[1]$.

\begin{exer} Show that $\sim$ satisfies
$$g\sim g'\implies hg\sim hg',\quad g\sim g'\implies gh\sim g'h$$
if and only if it satisfies
$$g\sim g',h\sim h'\implies gh\sim g'h'.$$
\end{exer}

So, to have a well-defined group structure on $G/\sim$, the left cosets of $[1]$ must be equal to the right cosets of $[1]$. That is, $[1]$ must be a normal subgroup. Tracing through all the previous proofs also shows that, for a normal $N\unlhd G$, quotienting by the equivalence relation $g\sim g'\iff gN=g'N$ can yield a well-defined group structure.

\begin{defn*} Let $G$ be a group and $N\unlhd G$ be a normal subgroup. The \textit{quotient group} $G/N$ is the set of left cosets of $N$ under the operation $(gN)(g'N)=(gg')N$.\end{defn*}

\begin{rem*} Let's show ``by hand" that this group operation is well-defined. Suppose $gN=g'N$ and $hN=h'N$. We want to show that $ghN=g'h'N$. We have that $g^{-1}g',h^{-1}h'\in N$. Now,
$$(gh)^{-1}(g'h')=h^{-1}(g^{-1}g')h'=h^{-1}nh',$$
where we write $n=g^{-1}g'\in N$. Since $N$ is normal, we can find $n'\in N$ with $nh'=h'n'$ (convince yourself that this follows from $(h')^{-1}Nh'\subseteq N$). Thus,
$$h^{-1}nh'=h^{-1}h'n'=(h^{-1}h')n'\in N,$$
as desired.
\end{rem*}

\begin{exer} We've shown that $G/N$ has a well-defined binary operation. Show that it is indeed a group with identity $1N$ and inverses $(gN)^{-1}=g^{-1}N$.\end{exer}

So why do we care about quotient groups? I'll give you two rough reasons. The first is the following theorem.

\begin{thm*}[First Isomorphism Theorem] Let $G$ be a group and $\phi:G\to H$ a homomorphism. Then $G/\ker\phi\cong\im\phi$.\end{thm*}

So quotient groups give us an intrinsic way of describing how homomorphic images of $G$ can look. It also in particular tells us that there are only so many ways that such images can look.

Here's our second:

\begin{thm*}[Fourth Isomorphism Theorem] Let $G$ be a group and $N\unlhd G$. Then there is an order-preserving bijection between subgroups of $G$ containing $N$ and subgroups of $G/N$.\end{thm*}

That is, the ``subgroup structure" of $G/N$ appears in the ``subgroup structure" of $G$. This is then useful because studying subgroups of $G/N$ tells us something about $G$. It also tends to be useful because we can then do things like induct on $|G|$ and use the inductive hypothesis on $G/N$.

\section{The $5/8$ths Theorem}

Let's go through an example of how to work with normal subgroups.

\begin{prop*} Let $G$ be a group. If $G/Z(G)$ is cyclic then $G$ is abelian.\end{prop*}
\begin{proof} Write $Z=Z(G)$. If $G/Z$ is cyclic, then it is generated by one element, say $G/Z=\langle gZ\rangle$. Take $h,k\in G$. Then $hZ\in\langle gZ\rangle$, so there is some $n\in\ZZ$ with $hZ=(gZ)^n=g^nZ$. This means we can find some $z_1\in Z$ with $h=g^nz_1$. Symmetrically, there are $m\in\ZZ$ and $z_2\in Z$ with $k=g^mz_2$.

Now, since $z_1,z_2\in Z$, they commute with all elements of $G$. So, we have that
$$hk=(g^nz_1)(g^mz_2)=g^n(z_1g^m)z_2\eqby{z_1\in Z}g^ng^mz_1z_2.$$
Since $g^ng^m=g^{n+m}=g^mg^n$, we have that this is equal to
$$g^mg^nz_1z_2\eqby{z_2\in Z}g^mz_2g^nz_1=kh.$$
That is, $hk=kh$, so $G$ is abelian. \end{proof}

In particular, since we know that the only group of order $p$ for $p$ a prime is the cyclic group $C_p$, we see that $|G/Z(G)|=[G:Z(G)]\neq p$ for all primes $p$ whenever $G$ is nonabelian. Since $G\neq Z(G)$ for nonabelian $G$, this also gives us the following corollary:

\begin{cor*} If $G$ is nonabelian, then $[G:Z(G)]\geq 4$.\end{cor*}

We'll use this as a lemma in the following.

\begin{thm*} Let $G$ be a finite group. If the probability that two randomly chosen elements of $G$ commute is $>5/8$, then $G$ is abelian.\end{thm*}
\begin{proof} Let $G$ be nonabelian. Let $x$ and $y$ be randomly chosen elements of $G$. Then,
$$\PP(x,y\text{ commute})=\PP\big(x,y\text{ comm}|x\in Z(G)\big)\cdot\PP(x\in Z(G))+\PP\big(x,y\text{ comm}|x\not\in Z(G)\big)\cdot\PP(x\not\in Z(G)).$$
Set $p=\PP(x\in Z(G))$. Since elements of $Z(G)$ by definition commute with all other elements of $G$, we have that
$$\PP(x,y\text{ comm})=1\cdot p+\PP\big(x,y\text{ comm}|x\not\in Z(G)\big)\cdot(1-p).$$
Now, $x$ and $y$ commuting is equivalent to $y\in Z_G(x)$. If $x\not\in Z(G)$, then $Z_G(x)\neq G$, and so $[G:Z_G(x)]\geq 2$ and $\PP(y\in Z_G(x))=\frac{|Z_G(x)|}{|G|}=\frac{1}{[G:Z_G(x)]}\leq\frac{1}{2}$. Thus,
$$\PP(x,y\text{ comm})=p+\PP\big(y\in Z_G(x)|x\not\in Z(G)\big)\cdot (1-p)\leq p+\frac{1}{2}(1-p)=\frac{1}{2}+\frac{1}{2}p.$$
Now, $p=\PP(x\in Z(G))=\frac{|Z(G)|}{|G|}=\frac{1}{[G:Z(G)]}$, which is $\leq\frac{1}{4}$ by the previous corollary. Thus,
$$\PP(x,y\text{ comm})=\frac{1}{2}+\frac{1}{2}p\leq \frac{1}{2}+\frac{1}{2}\cdot\frac{1}{4}=\frac{5}{8}.$$
\end{proof}

This bound is sharp: the group $Q_8=\{\pm 1, \pm i, \pm j, \pm k\}$ with multiplication $i^2=j^2=k^2=ijk=-1$ is not abelian but the above probability is exactly $5/8$.

\pagebreak

\section{Using Orbit-Stabilizer to Count}

\begin{prop*} Let $G$ be a finite group and $H<G$ be a proper subgroup. Then there exists some $k\in G$ with $k\not\in\bigcup_{g\in G}gHg^{-1}$. That is, there is some element in $G$ that is not conjugate to any element of $H$.\end{prop*}
\begin{proof} We'll first present the naive way of approaching this. We consider $G\acts\pset(G)$ via conjugation. The sets $gHg^{-1}$ are then the elements of the orbit of $H$. Since $|gHg^{-1}|=|H|$ (check this!), we have an upper bound
$$\left|\bigcup_{g\in G}gHg^{-1}\right|\leq |H|\cdot|\mc{O}_H|.$$
Now, $|\mc{O}_H|=[G:G_H]$ by orbit stabilizer. In this context, $G_H=N_G(H)$ is the normalizer, and so
$$\left|\bigcup_{g\in G}gHg^{-1}\right|\leq |H|\cdot[G:N_G(H)]=|G|\cdot\frac{|H|}{|N_G(H)|}.$$
If $|H|/|N_G(H)|<1$, this shows that we must be able to find some $g\in G$ not conjugate to any element of $H$. Unfortunately, this isn't always true.

To salvage this proof, we note that we were being overly conservative in our previous bound. We assumed that different conjugates of $H$ are disjoint. However, this isn't true: $1\in gHg^{-1}$ for all $g\in G$. So, we'll split $H$ into $\{1\}$ and $H\setminus\{1\}$. This gives us the sharper bound
$$\left|\bigcup_{g\in G}gHg^{-1}\right|=1+\big(|H|-1\big)\cdot|\mc{O}_H|=|G|\cdot\frac{|H|}{|N_G(H)|}-\big([G:N_G(H)]-1\big).$$
Just as before, if $|H|/|N_G(H)|<1$ then this is $<|G|$ and we are done. Suppose that $|H|=|N_G(H)|$. Since $H\leq N_G(H)$, this gives us that $N_G(H)=H<G$ is a proper subgroup of $G$. Thus, $[G:N_G(H)]\geq 2$, and so $[G:N_G(H)]-1\geq 1$. Thus, in either case we see that $\left|\bigcup_{g\in G}gHg^{-1}\right|<|G|$ and so there is some $k\in G$ not conjugate to any element of $H$. \end{proof}

\end{document}
